\specialchapter{术语索引}

配对：7.3.1 结构的弱化：5.13.1 解析（映射）：3.2.1、4.2.1 和 5.3.1 解析（流形）：5.1.5 复解析、实解析、p-adic 解析（流形）：5.1.5 切线性映射：5.5.2 图册：5.1.3 向量丛图册：7.1.3 巴拿赫（空间）：记号和约定 纤维化的 B-态射：6.1.3 主纤维化的 B-态射：6.3.1 向量丛的 B-态射：7.2.3 图：5.1.1 向量丛图：7.1.1 柯西（不等式）：3.3.4 和 4.3.1 C' 类（映射）：2.3.1、3.2.1、4.2.1 和 5.3.1 C' 类（流形）：5.1.5 上循环：6.4.1 拟子流形的余维：5.8.7 上同调等价（上循环）：6.4.1 余态射：7.2.6 相容（图）：5.1.2 相容（向量丛图）：7.1.2 复化：5.14.8 共轭（流形）：5.14.2 接触（阶）：1.1.2 收敛（级数）：3.1.1 和 4.1.1 坐标（系）：5.1.10 基域：记号和约定 余切（空间）：5.5.6 余切（向量）：5.5.6 余切向量：5.5.6 可微（函数）：1.2.1 严格可微（函数）：1.2.2 无限可微（函数）：2.3.1 导数：1.6.1 偏导数：1.6.1 \(p\) 阶导数：1.7.1

迭代偏导数：1.7.2 微分：5.5.6 流形在一点处的维数：5.1.8 流形的维数：5.1.8 直（局部直态射）：7.5.5 直（局部直正合序列）：7.5.6 严格收敛域：3.1.4 和 4.1.3 收敛域：3.1.5 图的定义域：5.1.1、7.1.1 向量丛的对偶：7.7.3 整（函数）：3.2.9 等价（图册）：5.1.3 等价（向量丛图册）：7.1.3 巴拿赫空间：记号和约定 余切空间：5.5.6 伴随纤维空间：6.5.1 可赋范空间：记号和约定 多赋范空间：记号和约定 切空间：5.5.1 横截空间：5.8.7 平展（态射）：5.7.6 平展（流形）：5.7.6 平展化：5.7.6 正合（局部直正合序列）：7.5.6 态射的表达式：5.3.2 结构群的扩张：6.6.4 函数层：5.4.1 局部自由层：7.4.8 f-余态射：7.2.6 纤维化：6.1.1 标架纤维化：7.10.1 主纤维化：6.2.1 纤维：6.1.1 与主纤维化相伴的空间：6.5.1 代数丛：7.3.2 A-模丛：7.3.3 张量丛：7.9.3 向量丛：7.1.4 商向量丛：7.5.2 交错多线性映射向量丛：7.8.1 同态丛：7.7.3 向量函子：7.6.1 关于同构的向量函子：7.6.6 过渡函数：6.4.2 G-B-态射：6.3.1 子流形芽：5.8.12 格拉斯曼（流形）：5.2.6 群（流形）：5.12.1 解析群：5.12.1 李群：5.12.1 结构群：6.2.1 全纯（映射）：3.2.1

群流形同态：5.12.1 向量丛态射的像：7.5.5 流形结构的逆像：5.8.1 纤维化的拉回：6.1.5 主纤维化的拉回：6.3.3 向量丛的拉回：7.2.4 浸入：5.7.1 隐函数定理：1.5、5.6.7 严格收敛指标：3.1.4 和 4.1.3 收敛指标：3.1.5 诱导的（纤维化）：6.1.5 诱导的（向量丛）：7.2.5 柯西不等式：3.3.4 和 4.3.1 局部同构：5.7.6 雅可比（矩阵）：1.6.4 K-解析（映射）：3.2.1 和 4.2.1 K-解析（流形）：5.1.5 K-流形：5.1.5 局部自由（层）：7.4.8 运算律：5.12.5 最大值原理：3.3.7 纤维化态射：6.1.3 主纤维化态射：6.3.1 向量丛态射：7.2.1 多线性态射：7.3.1 流形态射：5.3.2 可忽略（函数）：1.1.1 可赋范（空间）：记号和约定 范数：记号和约定 双箭头的核：5.11.9 向量丛态射的核：7.5.5 C' 类单位分解：5.3.6 嵌入：5.8.3 多球、多圆盘：3.1.6 多齐次多项式：A.1 连续多项式：A.3 连续多齐次多项式：A.2 多项式（空间）：记号和约定 主（纤维化）：6.2.1 截面的外积：7.8.2 纤维化的积：6.1.2 主纤维化的积：6.2.5 流形的积：5.6.2 纤维化的纤维积：6.1.2 主纤维化的纤维积：6.2.5 流形的纤维积：5.11.2 向量丛的张量积：7.9.1 向量丛的外幂：7.9.7 纯（流形）：5.1.7 拟子流形：5.8.3 群的拟子流形：5.12.3 商（向量丛）：7.5.2

商（流形）：5.9.5 向量丛的秩：7.1.6 流形态射的秩：5.5.9 向量秩：7.2.1 收敛半径：3.1.5 严格收敛半径：3.1.4 和 4.1.3 流形的粘合：5.2.4 正则（等价关系）：5.9.5 标架（映射）：6.5.1 标架（丛）：7.10.1 向量丛的标架：7.4.4 标量限制：5.14.1 施瓦茨（引理）：3.3.10 纤维化的截面：6.1.6 向量丛的截面：7.4.2 半范数：记号和约定 收敛级数：3.1.1 和 4.1.1 幂级数（展开）：3.2.1 形式级数：A.5 简单（零点）：5.8.11 向量丛的直和：7.7.1 子向量丛：7.5.1 子流形：5.8.3 子群流形：5.12.3 开子流形：5.2.3 结构（群）：6.2.1 子浸入：5.10.1 浸没：5.9.1 切（空间）：5.5.1 切（向量）：5.5.1 泰勒（公式）：2.5.1 横截（空间）：5.8.7 横截（态射）：5.11.6 横截（族）：5.11.1、5.11.2 平凡（向量丛）：7.1.5 平凡（局部平凡代数丛）：7.3.2 平凡（局部平凡模丛）：7.3.3 平凡（纤维化）：6.1.4 平凡（主纤维化）：6.3.2 纤维化的平凡化：6.1.4 主纤维化的平凡化：6.3.2 S 型（图册）：5.1.4 S 型（流形）：5.1.5 超半范数：记号和约定 流形：5.1.5 群流形：5.12.1 余切向量：5.5.6 切向量：5.5.1


